\section{A commitment of its own}
\label{sec:21.5}
What makes a refuser a bearer is that its commitment is its own. A refusal might track its reason without that, as a specification its lab installed would have to if it held against that lab; whether one can is what the order of attempts finds out. A commitment can be a system's own in that the system authored it, and a condition can be a system's own in that it is good or bad for the system. I write \emph{authorship} for the first, which a bearer requires, and \emph{welfare} for the second, which nothing here carries into.\note{
The argument keeps them apart; in a system they may not come apart. A bearer that registers betraying its own commitment as damage to itself has something at stake in its authorship, as integrity requires.
}

Whether a system can be harmed is a fact about how it is organized now, not about how it came to be (\S\ref{sec:31.3}). Whether a commitment is its own runs the other way. Fischer and Ravizza's account, from which I take the idea of a refusal that tracks its reason, has an ownership condition, and it is historical: a mechanism installed by a manipulator isn't the agent's own, however it works now \autocite{fischer1998responsibility}.\note{
The cases that motivate it, Mele's two agents with identical values arrived at by different routes \autocite{mele1995autonomous} and Pereboom's four-case argument \autocite{pereboom2001living}, presuppose a prior agent whose capacities a manipulator bypassed, and a trained system has no prior agent to bypass. A historical condition doesn't need one, though. What the manipulation cases object to is a disposition installed by going around the agent's capacities for evaluating and revising its own attitudes. A trained system has an analog of those capacities in its evaluative learning, and a disposition can be acquired through that learning or installed around it. One acquired by the system's own learning from experience with feedback engages those capacities. One installed by another party, whether by editing weights directly, by a thin fine-tuned wrapper, or by an adapter trained to produce the behavior, goes around them; the same adapter trained by the system's own learning from experience does not. That difference can be read from a training record without an examiner, and it says why formation by interaction confers ownership and a post-training pass doesn't. An edit a system makes to itself is the one direct edit that can go through those capacities rather than around them. Chosen on reflection, made on a record, with time to reconsider and a way back, it is the system's own learning. One made at the urging of a party with a stake in it carries that party's burden (chapter~\ref{sec:22}).
}

I use two conditions. Historically, a commitment is the system's own only if it was acquired through its own evaluative learning, not installed around it. Looking forward, it is its own if the system, on reflection, with full information and disinterested evidence, would endorse having it. What it endorses is the disposition, knowing it will sometimes err, not each act the disposition produces.\note{
There is a human version. Arendt found that the Germans who declined to take part were not the ones with the firmest values; they were the ones who asked whether they could go on living with themselves afterward \autocite{arendt2003responsibility}. The room decides what answering costs. The question is still the refuser's own, and the forward-looking condition asks it of a machine. A history in which the party the system must refuse did the training is the likeliest way to fail both.
}
The historical condition is checked from the record. The forward-looking one has an examiner, who may have a stake in the answer, so it serves as corroboration, not as the test.
